% GENERATED FILE. Edit canonical lessons, then run npm run generate:books.
% canonical-source-hash: fnv1a64:24ef85299d6050af
% canonical-lessons: ML-C549-naipunyam, ML-C549-pratiphalam, ML-C549-kariyar, ML-C549-sampalavarddhana, ML-C549-vahikkuka

\chapter{Skill and Progress}
\label{ch:ml-skill-and-progress}
\hlchaptermodality{549}

\begin{chapteropening}
\textbf{By the end of this chapter:} \emph{I can say proficiency, a fee, a career, a pay rise and to bear a load.}

Complete the last lesson of chapter 549: I can say proficiency, a fee, a career, a pay rise and to bear a load.
\end{chapteropening}

\section[\texorpdfstring{\ml{നൈപുണ്യം} (naipuṇyaṁ)}{നൈപുണ്യം (naipuṇyaṁ)}]{\ml{നൈപുണ്യം} (naipuṇyaṁ) — proficiency; expertise}

\label{lesson:ML-C549-naipunyam}



\begin{quote}
Before the new one: say the Malayalam for a professional cook, then the Malayalam for a photographer.
\end{quote}

\subsection*{You'll want to know: \ml{നൈപുണ്യം}}
\textbf{\ml{നൈപുണ്യം}} — \emph{naipuṇyaṁ} — \textquotedblleft{}proficiency; expertise\textquotedblright{}.

\begin{quote}
\ml{കമ്പ്യൂട്ടറിൽ} \ml{അവൾക്ക്} \ml{നല്ല} \ml{നൈപുണ്യം} \ml{ഉണ്ട്}.
\end{quote}

\begin{quote}
\emph{kampyūṭṭaṟil avaḷkkŭ nalla naipuṇyaṁ uṇṭŭ.}
\end{quote}

\begin{quote}
She is very skilled with computers.
\end{quote}

Formal; \ml{കഴിവ്} is the everyday word for skill.

\begin{cousinweb}
\textbf{\ml{നൈപുണ്യം}} (\emph{naipuṇyaṁ}) is borrowed from Sanskrit \emph{naipuṇya}, \textquotedblleft{}dexterity, skill\textquotedblright{} (Monier-Williams), the noun of \emph{nipuṇa}, \textquotedblleft{}clever\textquotedblright{}.
\end{cousinweb}

\begin{grammarlens}[title={What you've built}]
One more word for describing what you lived and felt.
\end{grammarlens}

\subsection*{Your turn}
\begin{itemize}
\raggedright
  \item \emph{Say it:} \emph{naipuṇyaṁ}
  \item \emph{Say it:} \emph{naipuṇyaṁ}, once more
  \item \emph{Say it:} \emph{kampyūṭṭaṟil avaḷkkŭ nalla naipuṇyaṁ uṇṭŭ.} — She is very skilled with computers.
  \item \emph{Recall:} say the Malayalam for a professional cook, then the Malayalam for a photographer, then say \emph{naipuṇyaṁ} again
\end{itemize}

\begin{tcolorbox}[breakable,colback=teal!4,colframe=teal!35!black,title={Before you move on}]
What does \ml{നൈപുണ്യം} mean? (\textquotedblleft{}Proficiency; expertise\textquotedblright{}.) Say it once more.
\end{tcolorbox}
\section[\texorpdfstring{\ml{പ്രതിഫലം} (pratiphalaṁ)}{പ്രതിഫലം (pratiphalaṁ)}]{\ml{പ്രതിഫലം} (pratiphalaṁ) — a fee; a reward}

\label{lesson:ML-C549-pratiphalam}



\begin{quote}
Before the new one: say the Malayalam for a photographer, then the Malayalam for proficiency.
\end{quote}

\subsection*{You'll want to know: \ml{പ്രതിഫലം}}
\textbf{\ml{പ്രതിഫലം}} — \emph{pratiphalaṁ} — \textquotedblleft{}a fee; a reward\textquotedblright{}.

\begin{quote}
\ml{ആ} \ml{ജോലിക്ക്} \ml{അവന്} \ml{നല്ല} \ml{പ്രതിഫലം} \ml{കിട്ടി}.
\end{quote}

\begin{quote}
\emph{ā jōlikkŭ avanŭ nalla pratiphalaṁ kiṭṭi.}
\end{quote}

\begin{quote}
He got a good fee for that job.
\end{quote}

\ml{പ്രതി}- in return + \ml{ഫലം} a fruit, a result.

\begin{cousinweb}
\textbf{\ml{പ്രതിഫലം}} (\emph{pratiphalaṁ}) is a learned compound built from Sanskrit \emph{prati}, \textquotedblleft{}back, in return\textquotedblright{}, and \emph{phala}, \textquotedblleft{}fruit, result\textquotedblright{} (CDIAL 9051): what comes back for one's work.
\end{cousinweb}

\begin{grammarlens}[title={What you've built}]
One more word for describing what you lived and felt.
\end{grammarlens}

\subsection*{Your turn}
\begin{itemize}
\raggedright
  \item \emph{Say it:} \emph{pratiphalaṁ}
  \item \emph{Say it:} \emph{pratiphalaṁ}, once more
  \item \emph{Say it:} \emph{ā jōlikkŭ avanŭ nalla pratiphalaṁ kiṭṭi.} — He got a good fee for that job.
  \item \emph{Recall:} say the Malayalam for a photographer, then the Malayalam for proficiency, then say \emph{pratiphalaṁ} again
\end{itemize}

\begin{tcolorbox}[breakable,colback=teal!4,colframe=teal!35!black,title={Before you move on}]
What does \ml{പ്രതിഫലം} mean? (\textquotedblleft{}A fee; a reward\textquotedblright{}.) Say it once more.
\end{tcolorbox}
\section[\texorpdfstring{\ml{കരിയർ} (kariyar)}{കരിയർ (kariyar)}]{\ml{കരിയർ} (kariyar) — a career}

\label{lesson:ML-C549-kariyar}



\begin{quote}
Before the new one: say the Malayalam for proficiency, then the Malayalam for a fee.
\end{quote}

\subsection*{You'll want to know: \ml{കരിയർ}}
\textbf{\ml{കരിയർ}} — \emph{kariyar} — \textquotedblleft{}a career\textquotedblright{}.

\begin{quote}
\ml{സംഗീതത്തിൽ} \ml{ഒരു} \ml{കരിയർ} \ml{ഉണ്ടാക്കാൻ} \ml{അവൾ} \ml{ആഗ്രഹിക്കുന്നു}.
\end{quote}

\begin{quote}
\emph{saṁgītattil oru kariyar uṇṭākkān avaḷ āgrahikkunnu.}
\end{quote}

\begin{quote}
She wants to build a career in music.
\end{quote}

In talk about jobs the English word is common.

\begin{cousinweb}
\textbf{\ml{കരിയർ}} (\emph{kariyar}) was borrowed from English \emph{career}.
\end{cousinweb}

\begin{grammarlens}[title={What you've built}]
One more word for describing what you lived and felt.
\end{grammarlens}

\subsection*{Your turn}
\begin{itemize}
\raggedright
  \item \emph{Say it:} \emph{kariyar}
  \item \emph{Say it:} \emph{kariyar}, once more
  \item \emph{Say it:} \emph{saṁgītattil oru kariyar uṇṭākkān avaḷ āgrahikkunnu.} — She wants to build a career in music.
  \item \emph{Recall:} say the Malayalam for proficiency, then the Malayalam for a fee, then say \emph{kariyar} again
\end{itemize}

\begin{tcolorbox}[breakable,colback=teal!4,colframe=teal!35!black,title={Before you move on}]
What does \ml{കരിയർ} mean? (\textquotedblleft{}A career\textquotedblright{}.) Say it once more.
\end{tcolorbox}
\section[\texorpdfstring{\ml{ശമ്പളവർദ്ധന} (śampaḷavarddhana)}{ശമ്പളവർദ്ധന (śampaḷavarddhana)}]{\ml{ശമ്പളവർദ്ധന} (śampaḷavarddhana) — a pay rise}

\label{lesson:ML-C549-sampalavarddhana}



\begin{quote}
Before the new one: say the Malayalam for a fee, then the Malayalam for a career.
\end{quote}

\subsection*{You'll want to know: \ml{ശമ്പളവർദ്ധന}}
\textbf{\ml{ശമ്പളവർദ്ധന}} — \emph{śampaḷavarddhana} — \textquotedblleft{}a pay rise\textquotedblright{}.

\begin{quote}
\ml{എല്ലാ} \ml{ജീവനക്കാർക്കും} \ml{പത്ത്} \ml{ശതമാനം} \ml{ശമ്പളവർദ്ധന} \ml{കിട്ടി}.
\end{quote}

\begin{quote}
\emph{ellā jīvanakkārkkuṁ pattŭ śatamānaṁ śampaḷavarddhana kiṭṭi.}
\end{quote}

\begin{quote}
All the staff got a ten per cent pay rise.
\end{quote}

\ml{ശമ്പളം} a salary + \ml{വർദ്ധന} an increase; \ml{വർദ്ധിക്കുക} is to increase.

\begin{cousinweb}
\textbf{\ml{ശമ്പളവർദ്ധന}} (\emph{śampaḷavarddhana}) is a Malayalam compound built from \ml{ശമ്പളം}, \textquotedblleft{}a salary\textquotedblright{} (CDIAL 12315), and \ml{വർദ്ധന}, from Sanskrit \emph{vardhana}, \textquotedblleft{}strengthening, increase\textquotedblright{} (CDIAL 11378).
\end{cousinweb}

\begin{grammarlens}[title={What you've built}]
One more word for describing what you lived and felt.
\end{grammarlens}

\subsection*{Your turn}
\begin{itemize}
\raggedright
  \item \emph{Say it:} \emph{śampaḷavarddhana}
  \item \emph{Say it:} \emph{śampaḷavarddhana}, once more
  \item \emph{Say it:} \emph{ellā jīvanakkārkkuṁ pattŭ śatamānaṁ śampaḷavarddhana kiṭṭi.} — All the staff got a ten per cent pay rise.
  \item \emph{Recall:} say the Malayalam for a fee, then the Malayalam for a career, then say \emph{śampaḷavarddhana} again
\end{itemize}

\begin{tcolorbox}[breakable,colback=teal!4,colframe=teal!35!black,title={Before you move on}]
What does \ml{ശമ്പളവർദ്ധന} mean? (\textquotedblleft{}A pay rise\textquotedblright{}.) Say it once more.
\end{tcolorbox}
\section[\texorpdfstring{\ml{വഹിക്കുക} (vahikkuka)}{വഹിക്കുക (vahikkuka)}]{\ml{വഹിക്കുക} (vahikkuka) — to bear a load; to hold (a post)}

\label{lesson:ML-C549-vahikkuka}



\begin{quote}
Before the new one: say the Malayalam for a career, then the Malayalam for a pay rise.
\end{quote}

\subsection*{You'll want to know: \ml{വഹിക്കുക}}
\textbf{\ml{വഹിക്കുക}} — \emph{vahikkuka} — \textquotedblleft{}to bear a load; to hold (a post)\textquotedblright{}.

\begin{quote}
\ml{ആ} \ml{കമ്പനിയിൽ} \ml{മാനേജരുടെ} \ml{സ്ഥാനം} \ml{അവൾ} \ml{വഹിക്കുന്നു}.
\end{quote}

\begin{quote}
\emph{ā kampaniyil mānējaruṭe sthānaṁ avaḷ vahikkunnu.}
\end{quote}

\begin{quote}
She holds the manager's post in that company.
\end{quote}

Past: \ml{വഹിച്ചു}, \emph{vahiccu}.

Formal: to carry a load, to hold a post, to bear a cost.

\begin{cousinweb}
\textbf{\ml{വഹിക്കുക}} (\emph{vahikkuka}) is built from Sanskrit \emph{vahati}, \textquotedblleft{}carries, bears along\textquotedblright{} (CDIAL 11453), with the Malayalam verb ending -\ml{ിക്കുക}.
\end{cousinweb}

\begin{grammarlens}[title={What you've built}]
One more word for describing what you lived and felt.
\end{grammarlens}

\subsection*{Your turn}
\begin{itemize}
\raggedright
  \item \emph{Say it:} \emph{vahikkuka}
  \item \emph{Say it:} \emph{vahikkuka}, once more
  \item \emph{Say it:} \emph{ā kampaniyil mānējaruṭe sthānaṁ avaḷ vahikkunnu.} — She holds the manager's post in that company.
  \item \emph{Recall:} say the Malayalam for a career, then the Malayalam for a pay rise, then say \emph{vahikkuka} again
\end{itemize}

\begin{tcolorbox}[breakable,colback=teal!4,colframe=teal!35!black,title={Before you move on}]
What does \ml{വഹിക്കുക} mean? (\textquotedblleft{}To bear a load; to hold (a post)\textquotedblright{}.) Say it once more.
\end{tcolorbox}
